% grip-methods-paper-v1.tex
% Paper 2: Novel GRIP Methods Paper
% Target venue: Journal of Graph Algorithms and Applications (JGAA)
\documentclass[11pt,twocolumn]{article}

\usepackage[margin=0.85in]{geometry}
\usepackage{amsmath,amssymb,amsthm,mathtools}
\usepackage[T1]{fontenc}
\usepackage{lmodern}
\usepackage{microtype}
\usepackage{graphicx}
\graphicspath{{../figures/}}
\usepackage{booktabs}
\usepackage{float}
\usepackage{dblfloatfix}  % fix figure*/table* placement in twocolumn
% Pseudocode environment (no algorithm.sty dependency)
\usepackage{fancyvrb}
\newcommand{\kw}[1]{\textbf{#1}}
\newcommand{\fn}[1]{\textsc{#1}}
\newcommand{\cmnt}[1]{\textit{\(\triangleright\) #1}}
\usepackage[hidelinks]{hyperref}
\usepackage{caption}
\usepackage{xcolor}
\usepackage{natbib}
\usepackage{url}

% Aggressive float placement to keep figures near their references
\renewcommand{\topfraction}{0.9}
\renewcommand{\bottomfraction}{0.9}
\renewcommand{\textfraction}{0.1}
\renewcommand{\floatpagefraction}{0.8}
\renewcommand{\dbltopfraction}{0.9}
\renewcommand{\dblfloatpagefraction}{0.8}
\setcounter{topnumber}{4}
\setcounter{bottomnumber}{4}
\setcounter{totalnumber}{8}
\setcounter{dbltopnumber}{4}

\newcommand{\RR}{\mathbb{R}}
\newcommand{\norm}[1]{\left\lVert #1 \right\rVert}
\newcommand{\vect}[1]{\mathbf{#1}}
\newcommand{\grip}{\texttt{grip}}
\newcommand{\bigO}{\mathcal{O}}

\title{Geometry-Aware Multiscale Graph Drawing\\with Weighted GRIP}

\author{
  Pawel Gajer\thanks{Center for Advanced Microbiome Research and Innovation
  (CAMRI), Institute for Genome Sciences, Department of Microbiology and
  Immunology, University of Maryland School of Medicine, Baltimore, MD, USA.
  \texttt{pgajer@gmail.com}}
  \and
  Jacques Ravel\thanks{Center for Advanced Microbiome Research and Innovation
  (CAMRI), Institute for Genome Sciences, Department of Microbiology and
  Immunology, University of Maryland School of Medicine, Baltimore, MD, USA.}
}

\date{}

\begin{document}
\maketitle

% ============================================================
% ABSTRACT
% ============================================================
\begin{abstract}
Classical GRIP is a multiscale graph drawing algorithm that builds a
hierarchy of vertex subsets through maximal independent set (MIS)
filtrations, places vertices coarse-to-fine, and refines with local
force-directed iterations at each level.  Because it uses combinatorial
graph distance (hop count) throughout its pipeline---in filtration
construction, neighborhood caching, insertion, and refinement---it cannot
preserve nontrivial edge-length geometry.  We present a family of
extensions that transform GRIP from a purely combinatorial multiscale
layout method into a geometry-aware multiscale family.  The main
contributions are: (1)~a weighted multiscale hierarchy based on weighted
MIS filtration construction using Dijkstra shortest-path radii;
(2)~weighted neighborhood caches, insertion anchors, and local
refinement forces driven by geodesic distances; (3)~coarse active-set
global repulsion that reduces multiscale foldovers in both the
combinatorial and weighted engines; (4)~integration of landmark geodesic
Kamada--Kawai (LGKK) refinement both within the multiscale hierarchy
(in-core) and as post-layout polish; and (5)~a systematic benchmark
framework over weighted synthetic graph families with 3D as the primary
evaluation space.  On a benchmark panel spanning meshes, tori, spheres,
Sierpinski carpets, porous 3D bodies, and irregular manifolds, weighted
GRIP with in-core LGKK refinement reduces mean full-geodesic-KK error by
22\% over combinatorial GRIP in 3D, while weighted GRIP with post-layout
LGKK polish achieves a 38\% reduction.  The KK-warm-started geodesic
polishing family (KK$\to$LGKK) remains the accuracy ceiling, but
weighted GRIP variants are substantially faster and require no external
KK initialization.  All methods are implemented in the open-source R
package \grip{} with a C++ core via Rcpp.
\end{abstract}

% ============================================================
% 1. INTRODUCTION
% ============================================================
\section{Introduction}\label{sec:intro}

Force-directed graph drawing is the most widely used family of methods
for producing readable graph layouts.  Classical algorithms such as
Fruchterman--Reingold (FR)~\citep{fruchterman1991} and
Kamada--Kawai (KK)~\citep{kamada1989} produce high-quality layouts for
small graphs, but their per-iteration cost of $\bigO(n^2)$ or
$\bigO(n^3)$ becomes prohibitive beyond a few hundred vertices.

Multiscale methods address this scalability barrier by constructing a
hierarchy of progressively coarser vertex subsets, placing vertices
coarse-to-fine, and refining with local force-directed iterations at each
level.  The GRIP algorithm~\citep{gajer2002,gajer2004} achieves this
through maximal independent set (MIS) filtrations, yielding near-linear
time complexity in practice and good layouts even for graphs with
thousands of vertices.  Related multiscale approaches include the work of
Harel and Koren~\citep{harel2002} and the multilevel force-directed
placement of Hu~\citep{hu2005}.

However, classical GRIP is fundamentally \emph{combinatorial}.  Every
operation in its pipeline---MIS filtration construction, neighborhood
caching, vertex insertion, and local refinement---uses hop-count graph
distance.  This design is natural for unweighted graphs, where the only
meaningful distance between vertices is the number of edges on a shortest
path.  But many graphs of practical interest carry nontrivial edge-length
geometry: mesh discretizations of curved surfaces, $k$-nearest-neighbor
graphs built from distance matrices, and biological networks with
weighted similarity edges all have edge lengths that encode geometric
information.  On such graphs, hop-count-based multiscale drawing cannot
preserve the intrinsic metric, because the hierarchy does not see it.

In this paper, we present a family of algorithmic extensions that
transform GRIP from a purely combinatorial multiscale layout method into
a \emph{geometry-aware multiscale family}.  The key idea is
straightforward: replace combinatorial operations with geodesic-distance
equivalents throughout the entire GRIP pipeline.  The resulting
\emph{weighted GRIP} algorithm uses weighted MIS filtrations, weighted
neighborhood caches, weighted insertion anchors, and weighted local
refinement forces.  We further integrate landmark geodesic Kamada--Kawai
(LGKK) refinement---a sparse geodesic stress minimizer---both within the
multiscale hierarchy and as a post-layout polish step.

We also describe improvements to the combinatorial engine that benefit
both the unweighted and weighted variants: coarse active-set global
repulsion to reduce multiscale foldovers, improved insertion-anchor
strategies, and alternative final-stage refinement modes.

Throughout this paper, we adopt 3D as the \emph{primary evaluation
space}.  The weighted graph families in our benchmark panel---meshes on
curved surfaces, tori, spheres, porous 3D bodies---are intrinsically
geometric objects that live naturally in $\RR^3$.  Forcing their layouts
into $\RR^2$ necessarily introduces a representation bottleneck.  We
retain 2D as an informative \emph{limitation track} that quantifies the
cost of dimensional reduction, but all primary conclusions are drawn from
the 3D results.

The methods described here are implemented in the open-source R package
\grip{}~\citep{gajer2002}, with a C++ computational core accessed via
Rcpp~\citep{eddelbuettel2011}.  A companion software
paper~\citep{gripRpaper} describes the package architecture, user-facing
APIs, and workflow tooling; the present paper focuses on the algorithmic
innovations and their empirical evaluation.


% ============================================================
% 2. BACKGROUND
% ============================================================
\section{Background}\label{sec:background}

\subsection{Classical GRIP}

The GRIP algorithm~\citep{gajer2002,gajer2004} builds a hierarchy of
vertex subsets through \emph{maximal independent set (MIS) filtrations}.
Given a connected graph $G = (V, E)$ with $|V| = n$, the filtration is a
nested sequence of vertex subsets
\[
  V = V_0 \supset V_1 \supset \cdots \supset V_L,
\]
where $V_L$ contains a small number of seed vertices and each $V_\ell$
is a maximal independent set of $G$ restricted to the vertices that
survived level~$\ell$.  The exclusion radius at level~$\ell$ is
$2^{\ell-1}$ hops: no two vertices in $V_\ell$ are within that
hop-count distance of each other.

Layout proceeds coarse-to-fine.  The seed vertices $V_L$ are placed
first (e.g., using random initialization or a small-graph layout).  At
each subsequent level~$\ell$, the vertices in $V_{\ell} \setminus
V_{\ell+1}$ are introduced one at a time.  Each newly introduced vertex
is placed near its nearest already-placed neighbors (insertion anchors)
and then refined with local Kamada--Kawai (KK) iterations that
adjust its position relative to a cached neighborhood of nearby vertices.

The neighborhood cache for each active vertex stores the $k$ nearest
vertices (by hop count) that are already placed.  The local KK
refinement uses these cached graph distances as target lengths, applying
spring forces to bring the Euclidean layout distances into agreement with
the graph-theoretic targets.

At the finest level ($\ell = 0$), all vertices are active and a final
round of refinement is applied.  In the original GRIP, this uses either
local KK or Fruchterman--Reingold (FR) style forces.

\subsection{Kamada--Kawai and Fruchterman--Reingold}

The Kamada--Kawai (KK) algorithm~\citep{kamada1989} minimizes a stress
energy that penalizes discrepancies between Euclidean layout distances
and graph-theoretic target distances:
\begin{equation}\label{eq:kk}
  E_{\mathrm{KK}}(X)
  = \frac{1}{2}\sum_{i < j}
    k_{ij}\bigl(\norm{\vect{x}_i - \vect{x}_j} - \ell_{ij}\bigr)^2,
\end{equation}
where $\ell_{ij} = L_0 g_{ij}$ is the target distance (proportional to
graph distance $g_{ij}$), $k_{ij} = K / g_{ij}^2$ is a stiffness
weight, and $L_0 = L/D$ scales the graph diameter~$D$ to a drawing
scale~$L$.

The Fruchterman--Reingold (FR) algorithm~\citep{fruchterman1991} uses
attractive spring forces along edges and repulsive forces between all
vertex pairs, with a simulated annealing temperature schedule.  FR does
not explicitly optimize a stress functional but produces layouts with
uniform edge lengths and good vertex separation.

\subsection{Geodesic distortion in weighted graphs}

When a graph $G = (V, E, w)$ carries positive edge weights
$w : E \to \RR_{> 0}$, the natural graph distance between vertices $i$
and $j$ is the shortest-path distance $g_{ij}$ computed as the minimum
total edge weight along any path from $i$ to~$j$.  For an embedding
$X = (\vect{x}_1, \ldots, \vect{x}_n)$ in $\RR^d$, we can measure the
\emph{embedding-induced geodesic distance} along a fixed graph shortest
path $\gamma_{ij}$ as the sum of Euclidean edge lengths along that path:
\begin{equation}\label{eq:geodesic}
  h_{ij}(X) = \sum_{e \in \gamma_{ij}} \norm{\vect{x}_{u(e)} - \vect{x}_{v(e)}}.
\end{equation}

The classical KK energy~\eqref{eq:kk} compares the graph distance
$g_{ij}$ with the straight-line chord distance
$\norm{\vect{x}_i - \vect{x}_j}$, which can be a poor proxy for
$h_{ij}(X)$ on non-planar or high-curvature graphs.  The
\emph{geodesic KK} (GKK) energy replaces the chord distance with the
embedding-induced geodesic distance:
\begin{equation}\label{eq:gkk}
  E_{\mathrm{geo}}(X)
  = \frac{1}{2}\sum_{i < j}
    k_{ij}\bigl(h_{ij}(X) - \ell_{ij}\bigr)^2.
\end{equation}

This full all-pairs geodesic functional is the most faithful measure of
layout quality for weighted graphs, but it is expensive to evaluate and
optimize for large graphs.

\subsection{Landmark geodesic KK (LGKK)}

The all-pairs GKK functional scales as $\bigO(n^2)$ in the number of
vertex pairs.  A practical sparse approximation retains two kinds of
vertex-pair interactions: (1)~local neighbors in the graph metric, and
(2)~a small set of far-away landmark vertices selected by farthest-point
sampling.  For each vertex~$i$, let $N(i)$ be its graph neighbors and
$L(i)$ a set of $M$ landmark vertices obtained by greedy farthest-point
sampling from~$i$.  The landmark-sparse geodesic KK (LGKK) functional
sums over the union of these local and long-range pairs:
\begin{equation}\label{eq:lgkk}
  E_{\mathrm{land}}(X)
  = \frac{1}{2}\sum_{(i,j) \in \mathcal{S}}
    k_{ij}\bigl(h_{ij}(X) - \ell_{ij}\bigr)^2,
\end{equation}
where $\mathcal{S} = \{(i,j) : j \in N(i) \cup L(i) \text{ or } i \in
N(j) \cup L(j)\}$.  This reduces the pair count from $\bigO(n^2)$ to
$\bigO(n(k + M))$ while preserving both local edge-scale information and
global shape constraints through the landmarks.


% ============================================================
% 3. LIMITATIONS OF CLASSICAL GRIP ON WEIGHTED GRAPHS
% ============================================================
\section{Limitations of Classical GRIP on Weighted Graphs}\label{sec:limitations}

Classical GRIP loses edge-length geometry at four distinct points in its
pipeline.  Understanding where the geometry is lost motivates the
weighted redesign.

\paragraph{Combinatorial MIS filtration.}
The filtration hierarchy uses hop-count exclusion radii: level~$\ell$
requires that no two retained vertices are within $2^{\ell-1}$ hops of
each other.  On a weighted graph whose edges vary substantially in
length, two vertices separated by many short edges may be metrically
close, while two vertices separated by a single long edge may be
metrically far.  The combinatorial filtration cannot distinguish these
cases, producing coarsening levels that are geometrically uneven.

\paragraph{Hop-count neighborhood caches.}
The local neighborhood cache for each active vertex stores the $k$
nearest vertices by hop count.  On a weighted graph, the $k$ nearest
vertices by hop count may not be the $k$ nearest vertices by geodesic
distance.  This means the local refinement forces at each level are
driven by the wrong set of neighbors.

\paragraph{Combinatorial insertion anchors.}
When a new vertex is introduced at level~$\ell$, its initial position is
computed from a small set of already-placed anchor vertices, selected as
the nearest retained vertices by hop count.  On weighted graphs, these
anchors may not be the geometrically nearest vertices, leading to poor
initial placement.

\paragraph{Hop-count refinement targets.}
The local KK refinement at each level uses graph distances as target
lengths.  In classical GRIP, these are hop-count distances, which do not
reflect the weighted metric.  The refinement therefore pulls vertices
toward positions that match the combinatorial structure but not the
geometric structure of the graph.

Together, these four points mean that classical GRIP treats all edges as
having equal length, regardless of their weights.  The result is a layout
that captures the \emph{topological} structure of the graph but not its
\emph{metric} structure.  For graphs with meaningful edge-length
geometry, this is a significant limitation.


% ============================================================
% 4. MODERN COMBINATORIAL GRIP IMPROVEMENTS
% ============================================================
\section{Modern Combinatorial GRIP Improvements}\label{sec:modern}

Before describing the weighted extensions, we present improvements to the
combinatorial engine that benefit both the unweighted and weighted
variants.

\subsection{Coarse active-set global repulsion}\label{sec:globalrep}

The classical GRIP refinement at coarse levels uses only local
KK forces computed from the cached neighborhood of each vertex.  Because
these neighborhoods are small, distant active regions at coarse levels
may drift toward each other and produce foldovers or self-overlaps before
the finer levels can correct them.

We address this by adding a weak global repulsion term on coarse MISF
levels ($\ell > 0$).  Let $V_\ell$ denote the active vertex set at
level~$\ell$ with $|V_\ell|$ active vertices.  For each active vertex
$v$, we compute an additional repulsive displacement from the other
active vertices:

\begin{itemize}
\item If $|V_\ell| \leq T_{\mathrm{exact}}$, repulsion is computed
  exactly against all other active vertices.
\item Otherwise, repulsion is approximated by sampling $s$ active
  vertices uniformly at random and scaling by $(|V_\ell| - 1)/s$.
\end{itemize}

The repulsive force between vertices $v$ and $u$ uses an inverse-square
form:
\[
  \vect{f}_{\mathrm{rep}}(v, u) =
  \alpha \cdot \frac{0.05 \cdot \bar{e}^2}{\norm{\vect{x}_v - \vect{x}_u}^2}
  \cdot \frac{\vect{x}_v - \vect{x}_u}{\norm{\vect{x}_v - \vect{x}_u}},
\]
where $\bar{e}$ is the mean edge length and $\alpha$ is the
\texttt{coarse\_repulsion\_factor} parameter.  This displacement is added
to the local KK displacement before temperature scaling.

The coarse global repulsion is disabled at the finest level ($\ell = 0$),
where the existing final-stage refinement (FR or KK-with-repulsion)
handles vertex separation.

\subsection{Improved insertion anchors}\label{sec:anchors}

Classical GRIP selects insertion anchors by a first-found BFS rule:
the first $K$ already-placed vertices encountered during BFS from the
newly introduced vertex become its placement anchors.  This is
order-sensitive and can produce geometrically skewed anchor sets.

We introduce three alternative anchor-selection strategies:

\begin{itemize}
\item \textbf{Distance-band}: expand BFS until the $K$-th anchor's
  distance band is exhausted, then select $K$ anchors from the full
  candidate pool.
\item \textbf{Balanced-band}: use the same band expansion, then
  explicitly select a subset whose centroid is centered in the candidate
  cloud while maintaining geometric spread.
\item \textbf{Spread-previous}: restrict anchors to the immediately
  previous MISF level and select for broad angular and geometric
  coverage.
\end{itemize}

These strategies reduce the order-sensitivity of insertion and improve
placement quality, especially on structured graph families where the BFS
visit order can be systematically biased.

\subsection{Alternative final-stage refinement}\label{sec:finalstage}

The original GRIP uses FR-style forces for the final all-vertices
refinement.  We add a \texttt{kk\_repulse} mode that replaces the FR
final stage with KK-style local distance-matching updates combined with
explicit active-set repulsion.  This mode maintains the stress-based
distance targets through the final stage rather than switching to the
purely spring-based FR model.

Additionally, we introduce a \texttt{final\_anchor\_factor} that
optionally anchors the final FR stage to the pre-final layout, preventing
excessive drift during the final all-vertices pass, and a
\texttt{final\_move\_scale\_after\_first} damping factor that reduces
displacement magnitudes after the first final-stage round.


% ============================================================
% 5. WEIGHTED GRIP
% ============================================================
\section{Weighted GRIP}\label{sec:weighted}

This section describes the core contribution: a geometry-aware weighted
counterpart to the combinatorial GRIP pipeline.  The weighted engine
replaces combinatorial operations with geodesic-distance equivalents at
every stage while preserving the same multiscale coarse-to-fine
structure.

\subsection{Edge-length normalization}\label{sec:normalization}

Before the weighted pipeline begins, edge lengths are globally normalized
so that the median edge weight becomes~1.  This ensures that the
weighted shortest-path radii used by the MISF hierarchy remain comparable
to the combinatorial radii $1, 2, 4, \ldots$ that govern the classical
filtration.  Mean normalization and no normalization are also supported.

\subsection{Weighted MIS filtration}\label{sec:wmisf}

The weighted MISF replaces the hop-count exclusion radius with a
weighted shortest-path radius.  At level~$\ell$, the exclusion radius is
$r_\ell = 2^{\ell-1}$ in the normalized weighted metric.  Two vertices
are excluded from co-existing at level~$\ell$ if their weighted
shortest-path distance (computed by Dijkstra in the full weighted graph)
is less than $r_\ell$.

This produces a coarsening hierarchy that respects the metric structure
of the graph: metrically close vertices are separated early, and the
spacing between retained vertices at each level is geometrically uniform
rather than topologically uniform.

The full procedure is given in Figure~\ref{alg:wmisf}.  The key
difference from the combinatorial algorithm is in line~8: the
neighborhood query uses Dijkstra shortest-path distance with a metric
radius cutoff, rather than BFS with a hop-count cutoff.  The radius
function $r(\ell) = 2^{\ell-1}$ produces the same geometric progression
as in classical GRIP, but measured in the normalized weighted metric.

\begin{figure}[t]
\centering
\fbox{\parbox{0.93\columnwidth}{\small
\textbf{Algorithm 1:} \fn{WeightedMISF}$(G, w, n_{\mathrm{init}})$\\[4pt]
\textbf{Input:} Weighted graph $G=(V,E,w)$; target coarsest-level
  size~$n_{\mathrm{init}}$\\
\textbf{Output:} MISF level assignment $\mathrm{level}[v]$ for each
  $v \in V$; level sizes $|V_0|, |V_1|, \ldots, |V_L|$\\[4pt]
\begin{tabular}{@{\hspace{0pt}}r@{\;\;}l}
1: & Normalize edge weights so $\mathrm{median}(w) = 1$ \\
2: & $\mathrm{level}[v] \gets 0$ for all $v \in V$;
     \; $\ell \gets 1$ \\
3: & $C \gets V$ \cmnt{candidate set for removal} \\
4: & \kw{while} $|C| > n_{\mathrm{init}}$ \kw{do} \\
5: & \quad $r_\ell \gets 2^{\ell-1}$
     \cmnt{weighted exclusion radius} \\
6: & \quad Shuffle $C$ randomly (Fisher--Yates) \\
7: & \quad \kw{for each} $v \in C$ in shuffled order \kw{do} \\
8: & \quad\quad $B_v \gets \{u \in C : d_w(v,u) \leq r_\ell\}$
     \cmnt{Dijkstra ball} \\
9: & \quad\quad \kw{if} no vertex in $B_v \setminus \{v\}$
     was already selected \kw{then} \\
10: & \quad\quad\quad Mark $v$ as \emph{selected} for level~$\ell$ \\
11: & \quad\quad\quad Mark all $u \in B_v \setminus \{v\}$ as
     \emph{removable} \\
12: & \quad\quad \kw{end if} \\
13: & \quad \kw{end for} \\
14: & \quad Move all removable vertices from $C$ to level~$\ell$:\\
   & \quad\quad $\mathrm{level}[u] \gets \ell$ for each removable $u$ \\
15: & \quad $\ell \gets \ell + 1$ \\
16: & \kw{end while} \\
17: & $L \gets \ell$; \; $V_L \gets C$
     \cmnt{coarsest seeds} \\
18: & \kw{return} $\mathrm{level}[\cdot]$, $\{|V_0|, \ldots, |V_L|\}$
\end{tabular}
}}
\caption{Weighted MIS filtration construction.  The critical difference
from classical GRIP is in line~8: the ball $B_v$ is defined by
Dijkstra weighted shortest-path distance $d_w$ with metric radius
$r_\ell$, rather than by BFS hop count.  Vertices are processed in
random order within each level to avoid systematic ordering bias.}
\label{alg:wmisf}
\end{figure}

\subsection{Weighted neighborhood caches}\label{sec:wcaches}

For each active vertex at level~$\ell$, the weighted engine caches the
$k$ nearest retained vertices by weighted shortest-path distance (rather
than hop count).  These caches are built by running Dijkstra from each
active vertex and collecting settled vertices in order of increasing
geodesic distance.

To reduce computational cost, the Dijkstra search uses
\emph{settled-order early stopping}: it terminates as soon as the
required neighborhood caches and insertion-anchor candidates are filled.
This is exact because the Dijkstra frontier is processed in
nondecreasing distance order.  Internal scratch arrays are reused across
searches to avoid repeated allocation.

For larger graphs, an optional \texttt{metric\_neighbor\_cap} parameter
limits the number of settled Dijkstra vertices per search, providing an
explicit approximation mode without changing the default exact semantics.

\subsection{Weighted insertion}\label{sec:winsertion}

When a new vertex $v$ is introduced at level~$\ell$, its initial position
is computed from weighted insertion anchors: the nearest already-placed
vertices by weighted shortest-path distance.  The same anchor-selection
strategies described in Section~\ref{sec:anchors} are available in
weighted form, with hop-count BFS replaced by Dijkstra traversal.

The initial position is computed by one of three rules:
\begin{itemize}
\item \textbf{Barycenter}: the centroid of the anchor positions, weighted
  by inverse geodesic distance.
\item \textbf{Circle} (2D only): placement on a circle centered at the
  weighted barycenter.
\item \textbf{Least squares}: a multi-anchor least-squares fit that
  minimizes the sum of squared discrepancies between Euclidean distances
  to anchors and the corresponding weighted graph distances.
\end{itemize}

After initial placement, a small number of local KK micro-polish steps
refine the position using weighted graph distances as targets.

The full weighted insertion procedure is given in
Figure~\ref{alg:winsertion}.  The anchor collection phase (lines~1--3)
uses Dijkstra traversal in place of BFS, and the initial-position
computation (lines~5--12) offers a least-squares mode that exploits
the weighted distance information for multi-anchor geometric fitting.

\begin{figure}[t]
\centering
\fbox{\parbox{0.93\columnwidth}{\small
\textbf{Algorithm 2:} \fn{WeightedInsert}$(v, \ell, X, G, w)$\\[4pt]
\textbf{Input:} New vertex~$v$ at MISF level~$\ell$; current layout~$X$;
  weighted graph $(G, w)$\\
\textbf{Output:} Initial position $\vect{x}_v \in \RR^d$\\[4pt]
\begin{tabular}{@{\hspace{0pt}}r@{\;\;}l}
\multicolumn{2}{@{}l}{\cmnt{Phase 1: Anchor collection via Dijkstra}} \\
1: & Run Dijkstra from $v$ in $(G, w)$ \\
2: & Collect anchors $A = \{a_1, \ldots, a_K\}$: the first $K$ \\
   & \quad already-placed vertices settled by Dijkstra, \\
   & \quad restricted to levels $> \ell$ \\
3: & Record $d_i = d_w(v, a_i)$ for each anchor $a_i$ \\[4pt]
\multicolumn{2}{@{}l}{\cmnt{Phase 2: Initial position}} \\
4: & \kw{if} placement mode = \texttt{least\_squares}
     and $K > d$ \kw{then} \\
5: & \quad Set up overdetermined system: for each $a_i$, \\
   & \quad\quad $\norm{\vect{x}_v - \vect{x}_{a_i}}^2 \approx
     (L_0 \cdot d_i)^2$ \\
6: & \quad Form normal equations
     $\vect{M}^{\!\top}\!\vect{M}\,\vect{x}_v
     = \vect{M}^{\!\top}\!\vect{b}$ \\
   & \quad\quad where row $i$: $\vect{m}_i = 2(\vect{x}_{a_i}
     - \vect{x}_{a_1})$, \\
   & \quad\quad $b_i = \norm{\vect{x}_{a_i}}^2
     - \norm{\vect{x}_{a_1}}^2
     - (L_0 d_i)^2 + (L_0 d_1)^2$ \\
7: & \quad Solve via Gaussian elimination with pivoting \\
8: & \quad \kw{if} system is singular \kw{then} fall back to
     barycenter \\
9: & \kw{else if} placement mode = \texttt{barycenter} \kw{then} \\
10: & \quad $\vect{x}_v \gets \sum_{i=1}^{K}
     \frac{d_i^{-1}}{\sum_j d_j^{-1}} \, \vect{x}_{a_i}$
     \cmnt{inverse-distance weighted} \\
11: & \kw{else} \cmnt{circle mode, 2D only} \\
12: & \quad Place $\vect{x}_v$ on circle through anchor centroid \\
13: & \kw{end if} \\[4pt]
\multicolumn{2}{@{}l}{\cmnt{Phase 3: Local KK micro-polish}} \\
14: & \kw{for} $t = 1, \ldots, t_{\mathrm{local}}$ \kw{do} \\
15: & \quad $\vect{\delta} \gets \sum_{i=1}^{K}
     \bigl(\frac{\norm{\vect{x}_v - \vect{x}_{a_i}}}{L_0 d_i}
     - 1\bigr)
     \cdot (\vect{x}_v - \vect{x}_{a_i})$ \\
16: & \quad $\vect{x}_v \gets \vect{x}_v - \eta\,\vect{\delta}$
     \cmnt{step with temperature $\eta$} \\
17: & \kw{end for} \\
18: & \kw{return} $\vect{x}_v$
\end{tabular}
}}
\caption{Weighted vertex insertion.  Phase~1 collects anchor vertices
using Dijkstra traversal (replacing BFS).  Phase~2 computes the initial
position; the least-squares mode solves for the position whose
Euclidean distances to anchors best match the scaled weighted
graph distances.  Phase~3 applies a few local KK micro-polish steps
using weighted distance targets.}
\label{alg:winsertion}
\end{figure}

\subsection{Weighted local refinement}\label{sec:wrefinement}

At each level~$\ell$, the weighted engine applies local KK-style
refinement using the weighted neighborhood caches.  The target distance
between vertex~$v$ and a cached neighbor~$u$ is the weighted
shortest-path distance $g_{vu}$ (in the normalized metric), and the
spring force pulls $v$ toward or away from $u$ to match this target.

The temperature schedule, adaptive momentum, and displacement clamping
follow the same structure as in the combinatorial engine, but the force
magnitudes are computed from geodesic distances rather than hop counts.

At the finest level ($\ell = 0$), the weighted engine supports the same
final-stage modes as the combinatorial engine: FR with weighted
edge-length targets, or KK-with-repulsion using weighted distances.

\subsection{Coarse global repulsion in weighted GRIP}

The coarse active-set global repulsion described in
Section~\ref{sec:globalrep} is equally available in the weighted engine.
The repulsive force computation is identical (it depends only on
Euclidean positions, not on graph distances), so the same parameters
control it in both engines.

\subsection{Algorithmic summary}

Table~\ref{tab:pipeline} summarizes the component-by-component
differences between the combinatorial and weighted GRIP pipelines, and
Figure~\ref{fig:pipeline} provides a visual flowchart.  The
two engines share the same coarse-to-fine multiscale structure and the
same parameter space; they differ in how distances are measured and how
neighborhoods are constructed.

\begin{table*}[!htb]
\centering
\caption{Component comparison between combinatorial GRIP and weighted
  GRIP.  All shared components use identical implementations; divergent
  components differ only in the distance metric used.}
\label{tab:pipeline}
\small
\begin{tabular}{lll}
\toprule
Component & Combinatorial GRIP & Weighted GRIP \\
\midrule
MISF exclusion radius & Hop count ($2^{\ell-1}$ hops) & Weighted shortest path ($2^{\ell-1}$ in normalized metric) \\
Distance computation & BFS & Dijkstra \\
Neighborhood caches & $k$ nearest by hop count & $k$ nearest by weighted shortest path \\
Insertion anchors & BFS-order selection & Dijkstra-order selection \\
Local KK targets & Hop-count distance & Weighted shortest-path distance \\
Final FR/KK-repulse & Hop-count neighborhoods & Weighted neighborhoods \\
Coarse global repulsion & Euclidean (shared) & Euclidean (shared) \\
Edge-length normalization & Not applicable & Median, mean, or none \\
\bottomrule
\end{tabular}
\end{table*}

\begin{figure*}[!htb]
\centering
\includegraphics[width=0.82\textwidth]{pipeline-comparison.pdf}
\caption{Side-by-side comparison of the combinatorial GRIP pipeline
  (left, blue) and the weighted GRIP pipeline (right, green).  Shared
  components---coarse global repulsion and optional LGKK
  polish---are shown in gray and purple respectively.  The weighted
  pipeline replaces BFS distances with Dijkstra distances throughout,
  adds an explicit weight-normalization step, and uses
  geodesic-aware radii for both MISF construction and neighborhood
  caching.}
\label{fig:pipeline}
\end{figure*}

% ============================================================
% 6. GEODESIC REFINEMENT
% ============================================================
\section{Geodesic Refinement Inside Weighted GRIP}\label{sec:lgkk}

The weighted GRIP engine described in Section~\ref{sec:weighted}
produces geometry-aware layouts, but its refinement is still
\emph{local}: each vertex is adjusted relative to a small neighborhood.
Landmark geodesic KK (LGKK) provides a complementary
\emph{semi-global} refinement that considers both local edge-scale
fidelity and long-range landmark constraints.

\subsection{Post-layout LGKK polish}

The simplest integration is to apply LGKK as a post-processing step
after the weighted GRIP solve completes.  The LGKK optimizer takes the
weighted GRIP layout as its starting point, constructs the sparse pair
set $\mathcal{S}$ (local neighbors plus farthest-point landmarks), and
runs gradient descent on the landmark-sparse geodesic
energy~\eqref{eq:lgkk} for a fixed number of iterations.

This is controlled by the \texttt{lgkk\_polish\_rounds} parameter.  The
key tuning knobs are the number of local neighbors
(\texttt{lgkk\_local\_nbrs}) and the number of landmarks per vertex
(\texttt{lgkk\_landmark\_count}).

\subsection{In-core multiscale LGKK refinement}

A more architecturally significant integration applies LGKK refinement
\emph{within} the multiscale hierarchy, after each eligible MISF level
completes its standard GRIP refinement rounds.  This allows the geodesic
correction to operate on progressively larger active sets as the
hierarchy unfolds, rather than waiting until the full graph is assembled.

The in-core LGKK integration supports staged budgets:
\begin{itemize}
\item \texttt{lgkk\_rounds\_coarse}: LGKK rounds on coarse MISF levels
  ($\ell > 1$),
\item \texttt{lgkk\_rounds\_pre\_final}: LGKK rounds on the last coarse
  level before the full graph opens ($\ell = 1$),
\item \texttt{lgkk\_rounds\_final}: LGKK rounds after the full-graph
  level completes ($\ell = 0$).
\end{itemize}

An \texttt{lgkk\_active\_limit} parameter skips the in-core LGKK stage
on levels where the active set exceeds a threshold, preventing expensive
geodesic cache construction on large intermediate levels.

\subsection{Relation to full GKK}

Full geodesic KK (GKK) optimizes the all-pairs geodesic
energy~\eqref{eq:gkk}, using every vertex pair.  It serves as the
accuracy reference in our benchmarks: the best achievable geodesic
fidelity for a given graph and dimension.  LGKK is its sparse
approximation, retaining local and landmark pairs only.  Weighted GRIP
is the scalable multiscale engine that provides fast initial layouts.

The recommended usage pattern is a cascade:
\begin{enumerate}
\item Weighted GRIP for fast geometry-aware initial layout,
\item In-core LGKK for incremental geodesic correction during the
  multiscale solve,
\item Optional post-layout LGKK polish for additional refinement.
\end{enumerate}

For benchmarking, we also evaluate a KK-warm-started geodesic polish
cascade (KK $\to$ GKK or KK $\to$ LGKK), which provides the
highest-accuracy reference at the cost of requiring an external KK
initialization and substantially higher runtime.


% ============================================================
% 7. EXPERIMENTAL DESIGN
% ============================================================
\section{Experimental Design}\label{sec:experiments}

\subsection{Graph families}

We evaluate on a panel of weighted synthetic graph families that span a
range of geometric complexities.  Each family is generated by the
\grip{} package's graph-family constructors, which produce graphs with
known ground-truth embeddings.  The families are:

\begin{enumerate}
\item \textbf{Mesh} (saddle surface): $6 \times 6$ rectangular grid with
  Euclidean edge weights induced by a saddle surface
  ($z = x^2 - y^2$).  36 vertices, 60 edges.
\item \textbf{Torus} (pinched): $6 \times 6$ toroidal grid with
  curvature-varying edge weights.  36 vertices, 72 edges.
\item \textbf{Sierpinski carpet} (ripple): level-3 Sierpinski carpet
  with ripple-surface edge weights.  64 vertices, 88 edges.
\item \textbf{Cube channel network} (twisted): 3D porous body with
  twisted channel geometry.  104 vertices, 212 edges.
\item \textbf{Irregular torus} (pinched): irregular triangulated torus
  with pinch deformation.  60 vertices, 188 edges.
\end{enumerate}

These families cover regular grids (mesh), closed manifolds (torus),
fractal structures (carpet), porous 3D bodies (cube channel), and
irregular manifolds (irregular torus).  Edge-weight coefficients of
variation range from 0.03 (near-uniform) to 0.55 (highly variable),
ensuring that the benchmark tests geometry awareness across a range of
metric complexities.

\subsection{Methods}

We compare seven methods organized into two families:

\paragraph{GRIP family:}
\begin{itemize}
\item \textbf{GRIP}: combinatorial GRIP with coarse global repulsion
  (the current default engine).
\item \textbf{Weighted GRIP}: weighted GRIP with coarse global repulsion,
  no LGKK refinement.
\item \textbf{Weighted GRIP + core LGKK}: weighted GRIP with in-core
  multiscale LGKK refinement.
\item \textbf{Weighted GRIP + polish LGKK}: weighted GRIP with
  post-layout LGKK polish.
\end{itemize}

\paragraph{KK family:}
\begin{itemize}
\item \textbf{KK}: Kamada--Kawai layout from igraph~\citep{csardi2006}
  with edge weights.
\item \textbf{KK $\to$ GKK}: KK layout polished by full geodesic KK
  optimization (12 iterations, profiled scaling).
\item \textbf{KK $\to$ LGKK}: KK layout polished by landmark geodesic
  KK optimization (12 iterations, 8 local neighbors, 10 landmarks).
\end{itemize}

All KK-family methods share a common warm-start: a crude initial layout
constructed from the flat parameter coordinates with small Gaussian noise,
followed by weighted KK optimization.  This ensures that the KK baseline
is already well-initialized before geodesic polishing begins.

\subsection{Dimensions}

Each case is evaluated in both $d = 3$ (primary track) and $d = 2$
(limitation track).  The 3D track is primary because the benchmark
families are intrinsically geometric objects; the 2D track quantifies the
cost of forcing an intrinsically 3D metric into two dimensions.

\subsection{Metrics}\label{sec:metrics}

The primary quality metric is the \textbf{full GKK relative RMSE}: the
root-mean-square relative error of the all-pairs geodesic KK energy,
computed by \texttt{grip.score.geodesic.kk()} with profiled optimal
scaling.  This measures how faithfully the embedding preserves the
weighted graph metric as measured along embedded graph paths.

We also report:
\begin{itemize}
\item \textbf{Classical KK relative RMSE}: the chord-based KK error,
  measuring fidelity to the simpler straight-line distance model.
\item \textbf{Procrustes RMSE}: alignment error against the ground-truth
  surface embedding (3D) or parameter coordinates (2D).
\item \textbf{Runtime}: wall-clock elapsed time in seconds.
\end{itemize}

\subsection{Seeds and reproducibility}

All methods use deterministic seeding.  The GRIP-family methods use the
\texttt{seed} parameter of the \grip{} package.  The KK-family methods
use the same shared initial layout for each graph and dimension.
Benchmark scripts and data are included in the package repository.


% ============================================================
% 8. RESULTS
% ============================================================
\section{Results}\label{sec:results}

We organize results around five questions that characterize the
contribution of each algorithmic component.

\subsection{Does weighted GRIP improve geodesic fidelity over
combinatorial GRIP?}

Table~\ref{tab:overall-d3} shows the aggregate 3D results.  Weighted
GRIP alone achieves a mean full-GKK relative RMSE of 0.231, compared
with 0.189 for combinatorial GRIP.  This apparent increase reflects a
tradeoff: weighted GRIP changes the optimization landscape by introducing
geodesic distance targets, which can increase the geodesic error on some
families where the combinatorial engine happens to find a configuration
that scores well under the GKK metric by chance.

However, when weighted GRIP is combined with LGKK refinement, the
picture changes substantially.  Weighted GRIP with in-core LGKK achieves
0.146, a 22.6\% reduction over combinatorial GRIP.  Weighted GRIP with
post-layout LGKK polish reaches 0.118, a 37.7\% reduction.  The
improvement is consistent across families
(Table~\ref{tab:perfamily-d3-full}).

\begin{table}[t]
\centering
\caption{Overall means for $d = 3$ across the benchmark panel.}
\label{tab:overall-d3}
\small
\begin{tabular}{lrr}
\toprule
Method & GKK rel.\ RMSE & Runtime (s) \\
\midrule
GRIP                       & 0.189 & 0.042 \\
Weighted GRIP              & 0.231 & 0.047 \\
W.\ GRIP + core LGKK      & 0.146 & 0.055 \\
W.\ GRIP + polish LGKK    & 0.118 & 0.216 \\
KK                         & 0.045 & 0.003 \\
KK $\to$ GKK              & 0.028 & 0.956 \\
KK $\to$ LGKK             & 0.023 & 0.153 \\
\bottomrule
\end{tabular}
\end{table}

\subsection{How much does coarse global repulsion help?}

The coarse global repulsion is embedded in the ``GRIP'' baseline
throughout our experiments (it is the current default engine).  Its
effect was validated during development on structured families prone to
multiscale foldovers.  In the mesh benchmark suite (3D, three mesh
geometries at three sizes each), the global-repulsion engine reduced the
incidence of foldover artifacts compared with the legacy engine
(\texttt{grip.layout.legacy()}), particularly on the larger $12 \times
12$ meshes where the coarsest MISF levels have enough active vertices for
distant regions to interact.

The benefit is structural rather than metric: global repulsion primarily
prevents topological defects (folds, overlaps) that would otherwise
require post-hoc correction, rather than directly improving the geodesic
stress.

\subsection{How much does LGKK refinement add to weighted GRIP?}

Comparing the weighted GRIP variants in 3D (Table~\ref{tab:overall-d3}):

\begin{itemize}
\item In-core LGKK reduces the mean GKK relative RMSE from 0.231
  (weighted GRIP alone) to 0.146, a 36.8\% improvement.
\item Post-layout LGKK polish reduces it further to 0.118, a 49.0\%
  improvement over weighted GRIP alone.
\end{itemize}

The in-core variant is architecturally more significant because it
incorporates geodesic correction during the multiscale hierarchy rather
than depending solely on terminal cleanup.  The post-layout variant
achieves marginally better accuracy on this panel but at approximately
$4\times$ the runtime of the in-core variant.

\subsection{How close does weighted GRIP get to the KK family?}

The KK-warm-started geodesic polishing family remains the accuracy
ceiling.  In 3D, KK $\to$ LGKK achieves a mean GKK relative RMSE of
0.023, compared with 0.118 for weighted GRIP + polish LGKK.  However,
the KK family requires an external KK initialization step that itself
costs nontrivial time on larger graphs, and the full GKK polish
(KK $\to$ GKK) is roughly $4.4\times$ slower than the LGKK variant.

The per-family results (Table~\ref{tab:perfamily-d3-full}) show that the gap
between the GRIP and KK families varies by graph structure.  On the mesh
family, which has nearly uniform edge weights and smooth geometry, the
gap is largest.  On the irregular torus and carpet families, which have
higher geometric complexity, the weighted GRIP + LGKK variants close
more of the gap.

Table~\ref{tab:perfamily-d3-full} provides the full per-family
breakdown, and Figure~\ref{fig:perfamily-bars} visualizes the
per-family comparison across all seven methods.

\begin{table*}[!htb]
\centering
\caption{Full-GKK relative RMSE by method and family in $d = 3$.  Lower
  is better.  The best result in each family is shown in bold.  The
  ``GRIP-family best'' column gives the best result among the four
  GRIP-family methods.}
\label{tab:perfamily-d3-full}
\small
\begin{tabular}{lrrrrrrr}
\toprule
 & \multicolumn{4}{c}{GRIP family} & \multicolumn{3}{c}{KK family} \\
\cmidrule(lr){2-5} \cmidrule(lr){6-8}
Family & GRIP & W.\ GRIP & + core & + polish
       & KK & $\to$GKK & $\to$LGKK \\
\midrule
Mesh saddle     & 0.095 & 0.099 & 0.030 & 0.015
                & 0.012 & \textbf{0.002} & 0.006 \\
Torus pinched   & 0.142 & 0.168 & 0.130 & 0.121
                & 0.074 & 0.041 & \textbf{0.039} \\
Carpet ripple   & 0.172 & 0.174 & 0.161 & 0.135
                & 0.058 & 0.041 & \textbf{0.031} \\
Cube channel    & 0.072 & 0.079 & 0.038 & 0.016
                & 0.036 & 0.011 & \textbf{0.006} \\
Irreg.\ torus  & 0.462 & 0.635 & 0.373 & 0.302
                & 0.044 & 0.045 & \textbf{0.032} \\
\bottomrule
\end{tabular}
\end{table*}

\begin{figure*}[!htb]
\centering
\includegraphics[width=0.92\textwidth]{fig4_gkk_rmse_bars.pdf}
\caption{Full-GKK relative RMSE by graph family in $d = 3$, for all
  seven methods.  Lower bars indicate better geodesic fidelity.  The
  irregular torus is by far the most challenging family, with GRIP-family
  methods remaining above 0.12 even after LGKK polish, while KK-family
  methods reach below~0.05.  On the mesh family, weighted GRIP + polish
  LGKK approaches the KK baselines.}
\label{fig:perfamily-bars}
\end{figure*}

\subsection{Per-family analysis}\label{sec:perfamily}

The five benchmark families exhibit qualitatively different responses to
the weighted GRIP pipeline, and these differences illuminate the
strengths and limitations of the approach.

\paragraph{Mesh saddle.}
The $6 \times 6$ saddle mesh has nearly uniform edge weights
(edge-weight CV\,=\,0.03 after normalization) and smooth, low-curvature
geometry.  This is the easiest family for all methods.  Pure KK already
achieves 0.012 GKK relative RMSE, and KK\,$\to$\,GKK reduces it to
0.002.  The GRIP family starts at 0.095 (combinatorial) and 0.099
(weighted).  The near-identical scores for these two reflect the
near-uniform weights: when all edges have similar length, weighted
coarsening behaves similarly to hop-count coarsening.  LGKK refinement
is effective here: in-core LGKK reduces the error to 0.030, and
post-layout polish reaches 0.015---still $7.5\times$ the KK\,$\to$\,GKK
result, but an $84\%$ reduction from raw GRIP.  The mesh is the family
where the gap between multiscale and direct optimization is widest in
relative terms, because the simple geometry allows KK to find a
near-optimal solution cheaply.

\paragraph{Torus pinched.}
The closed toroidal manifold with curvature-varying (``pinched'') edge
weights (CV\,=\,0.54) is substantially harder.  Combinatorial GRIP
scores 0.142, and weighted GRIP is \emph{worse} at 0.168.  This is a
case where the weighted coarsening changes the optimization landscape
without improving the outcome---the weighted MISF hierarchy separates
the pinched region differently, but the resulting initial layout is not
closer to the geodesic optimum.  However, LGKK refinement recovers
the situation: in-core LGKK brings the error to 0.130, and post-layout
polish reaches 0.121.  The KK family achieves 0.039 (KK\,$\to$\,LGKK),
about $3\times$ better.  The torus illustrates that for closed manifolds
with high curvature variation, the multiscale engine provides a
reasonable starting point but geodesic refinement is essential.

\paragraph{Sierpinski carpet ripple.}
The fractal carpet with ripple-surface weights (CV\,=\,0.48) is the
most resistant family to \emph{all} methods.  Even KK\,$\to$\,LGKK only
reaches 0.031.  Within the GRIP family, weighted GRIP (0.174) barely
differs from combinatorial GRIP (0.172), suggesting that the fractal
connectivity---not the edge-length variation---is the dominant source of
difficulty.  The recursive hole structure creates long-range constraints
that local multiscale refinement cannot easily resolve.  Post-layout
LGKK polish reduces the error to 0.135, a 22\% improvement but still
$4.3\times$ the KK-family ceiling.  This family highlights that
fractal graph topology remains a challenge for multiscale methods
regardless of whether they are geometry-aware.

\paragraph{Cube channel network.}
The 3D porous body with twisted channels (CV\,=\,0.03) is a graph
where the multiscale approach works relatively well.  Combinatorial GRIP
achieves 0.072, among the best GRIP-family scores across families.  The
3D porous structure, with its relatively uniform edge weights and
localized connectivity, suits multiscale coarsening.  Weighted GRIP is
slightly worse at 0.079, but in-core LGKK brings it to 0.038, and
post-layout polish reaches 0.016---only $2.8\times$ the
KK\,$\to$\,LGKK result of 0.006.  This is the family where the
weighted GRIP + LGKK pipeline comes closest to the KK-family accuracy in
absolute terms.

\paragraph{Irregular torus.}
The irregular triangulated torus with pinch deformation (CV\,=\,0.55) is
the hardest family for the GRIP pipeline.  Combinatorial GRIP scores
0.462, and weighted GRIP is dramatically worse at 0.635.  The irregular
triangulation creates a highly non-uniform mesh where the weighted MISF
hierarchy struggles to find geometrically meaningful coarsening levels.
LGKK refinement provides large improvements---in-core LGKK reduces
the error to 0.373 (a 41\% drop from weighted GRIP), and post-layout
polish reaches 0.302 (a 52\% drop)---but the residual error is still
$9.3\times$ the KK\,$\to$\,LGKK result of 0.032.  This family
represents the current frontier of difficulty for geometry-aware
multiscale drawing: the irregular connectivity and high weight variation
demand the kind of global optimization that multiscale methods are
specifically designed to approximate.  On such graphs, the GRIP pipeline
is best understood as a fast initializer for subsequent geodesic
refinement rather than as a standalone solution.

\subsection{What are the runtime tradeoffs?}

\begin{figure}[H]
\centering
\includegraphics[width=\columnwidth]{fig5_pareto_scatter.pdf}
\caption{Runtime--quality Pareto scatter in $d = 3$.  Each point
  represents one (method, family) pair.  Color encodes method; marker
  shape encodes graph family.  Methods closer to the lower-left corner
  offer the best quality-for-time trade-off.  The GRIP-family methods
  cluster at low runtimes but moderate quality; the KK $\to$ GKK
  pipeline yields the best quality at higher runtime cost; KK $\to$ LGKK
  and weighted GRIP + polish LGKK provide intermediate trade-offs.}
\label{fig:pareto}
\end{figure}

Figure~\ref{fig:pareto} shows the runtime--quality trade-off as a
scatter plot, and Table~\ref{tab:overall-d3} gives the aggregate
numbers.
Combinatorial GRIP and weighted GRIP are the fastest methods (0.042s and
0.047s respectively), reflecting the near-linear scaling of the
multiscale engine.  Adding in-core LGKK increases runtime modestly to
0.055s.  Post-layout LGKK polish is more expensive at 0.216s, still
substantially cheaper than the KK $\to$ GKK pipeline at 0.956s.

The KK $\to$ LGKK combination (0.153s) is faster than weighted GRIP +
polish LGKK (0.216s) on these small graphs, but this comparison is
dominated by the KK initialization step, which scales as
$\bigO(n^2 \cdot \text{iterations})$.  On larger graphs, the multiscale
GRIP approach is expected to scale more favorably.

\subsection{2D limitation track}

Table~\ref{tab:overall-d2} shows the aggregate 2D results.  The same
method ranking persists at the aggregate level---KK $\to$ LGKK achieves
the lowest mean geodesic error, followed by KK $\to$ GKK, then the
weighted GRIP + LGKK variants---but the absolute error levels are
substantially higher.  This confirms the design decision to treat 3D as
the primary evaluation space: the weighted families in our benchmark
panel are intrinsically geometric objects, and 2D layouts systematically
face a representation bottleneck.

\begin{table}[t]
\centering
\caption{Overall means for $d = 2$ across the benchmark panel.}
\label{tab:overall-d2}
\small
\begin{tabular}{lrr}
\toprule
Method & GKK rel.\ RMSE & Runtime (s) \\
\midrule
GRIP                       & 0.185 & 0.033 \\
Weighted GRIP              & 0.246 & 0.054 \\
W.\ GRIP + core LGKK      & 0.152 & 0.059 \\
W.\ GRIP + polish LGKK    & 0.112 & 0.225 \\
KK                         & 0.089 & 0.002 \\
KK $\to$ GKK              & 0.066 & 1.007 \\
KK $\to$ LGKK             & 0.054 & 0.158 \\
\bottomrule
\end{tabular}
\end{table}

\subsection{Per-family dimensional comparison}\label{sec:dimcompare}

\begin{figure*}[!htb]
\centering
\includegraphics[width=0.92\textwidth]{fig7_2d_vs_3d.pdf}
\caption{2D vs.\ 3D geodesic fidelity for three representative methods:
  weighted GRIP + polish LGKK (the best GRIP-family variant), KK $\to$
  LGKK, and KK $\to$ GKK.  Hatched bars show 2D results; solid bars
  show 3D results.  Moving from 3D to 2D uniformly degrades quality,
  but the degradation is most pronounced on families with complex
  topology (torus, irregular torus) and less severe on the mesh.}
\label{fig:dimcompare}
\end{figure*}

Table~\ref{tab:2dvs3d} compares the best GRIP-family and best KK-family
results across dimensions for each graph family.  The comparison reveals
two noteworthy patterns beyond the expected overall 3D advantage.

\begin{table*}[!htb]
\centering
\caption{Best full-GKK relative RMSE by family, comparing $d = 2$ and
  $d = 3$.  For each family, the best GRIP-family method (always
  W.~GRIP + polish LGKK) and the best KK-family method are shown.
  Families where the GRIP-family method \emph{beats} the KK-family
  method in 2D are marked with~$\star$.}
\label{tab:2dvs3d}
\small
\begin{tabular}{l rr c rr c}
\toprule
 & \multicolumn{2}{c}{Best GRIP-family} & &
   \multicolumn{2}{c}{Best KK-family} & \\
\cmidrule(lr){2-3} \cmidrule(lr){5-6}
Family & $d{=}2$ & $d{=}3$ & $\Delta$ &
         $d{=}2$ & $d{=}3$ & $\Delta$ \\
\midrule
Mesh saddle     & 0.015 & 0.015 & $0\%$
                & 0.002 & 0.002 & $0\%$ \\
Torus pinched${}^{\star}$
                & 0.054 & 0.121 & $-55\%$
                & 0.102 & 0.039 & $+161\%$ \\
Carpet ripple   & 0.142 & 0.135 & $+6\%$
                & 0.053 & 0.031 & $+71\%$ \\
Cube channel${}^{\star}$
                & 0.027 & 0.016 & $+72\%$
                & 0.035 & 0.006 & $+517\%$ \\
Irreg.\ torus  & 0.323 & 0.302 & $+7\%$
                & 0.072 & 0.032 & $+125\%$ \\
\bottomrule
\end{tabular}
\end{table*}

First, the KK-family methods are \emph{far more sensitive to
dimensionality} than the GRIP-family methods.  On the cube channel
network, the best KK-family error increases by $517\%$ when moving from
3D to 2D (0.006 $\to$ 0.035), while the best GRIP-family error
increases by only $72\%$ (0.016 $\to$ 0.027).  On the torus, the
KK-family error increases by $161\%$ (0.039 $\to$ 0.102), while the
GRIP-family error actually \emph{decreases} by $55\%$
(0.121 $\to$ 0.054).

Second, on two families---the torus and the cube channel
network---weighted GRIP + polish LGKK \emph{beats} KK $\to$ LGKK in
2D.  On the pinched torus, the GRIP-family achieves 0.054 versus 0.102
for the KK-family.  On the cube channel, the GRIP-family achieves 0.027
versus 0.035.  These reversals do not occur in 3D, where the KK family
consistently dominates.

The explanation is that when forced into an insufficient ambient
dimension, the KK-based chord-distance optimization can be misled:
straight-line distances in $\RR^2$ become poor proxies for geodesic
distances on an intrinsically 3D surface, and the KK optimizer can
converge to a configuration that looks good under the chord metric but
scores poorly under the geodesic metric.  The multiscale GRIP engine,
which builds the layout incrementally through its coarsening hierarchy,
appears to be less susceptible to this particular failure mode in reduced
dimension.

This observation does not change the primary recommendation to use 3D
for these graph families; it does suggest that when 2D is unavoidable,
the weighted GRIP + LGKK pipeline deserves consideration as an
alternative to the KK-warm-start pathway.

\subsection{Mesh benchmark: KK/GKK/LGKK baseline}

A separate mesh-focused benchmark on three mesh geometries (saddle,
paraboloid, ripple) at three sizes ($8 \times 8$, $10 \times 10$,
$12 \times 12$) provides a more detailed picture of the geodesic
polishing baselines.  In 3D:

\begin{itemize}
\item Mean full-GKK relative RMSE dropped from 0.025 (pure KK) to 0.007
  (KK $\to$ GKK), a 73.7\% reduction.
\item KK $\to$ LGKK achieved 0.007, capturing nearly all of the GKK
  gain while being approximately $7\times$ faster.
\item The geodesic improvements came with only a modest 3--4\% increase
  in classical KK error, a mild tradeoff.
\end{itemize}

These results establish that LGKK is a strong practical approximation of
full GKK, and that the small accuracy gap between them does not justify
the much higher cost of full GKK for most applications.


% ============================================================
% 9. DISCUSSION
% ============================================================
\section{Discussion}\label{sec:discussion}

\subsection{When to use which variant}

The results suggest a clear decision framework:

\begin{itemize}
\item For \textbf{unweighted graphs} or graphs where edge weights do not
  encode geometry, use \textbf{combinatorial GRIP} (the default engine).
  It is the fastest method and handles topological structure well.
\item For \textbf{weighted graphs with meaningful edge lengths} where
  speed is the primary concern, use \textbf{weighted GRIP} alone.  It
  provides geometry-aware coarsening and refinement at minimal additional
  cost.
\item For \textbf{higher geodesic fidelity} on weighted graphs, use
  \textbf{weighted GRIP with in-core LGKK}.  This provides the best
  balance of speed and geodesic accuracy within the multiscale framework.
\item For \textbf{maximum geodesic accuracy} without regard to runtime,
  use the \textbf{KK $\to$ LGKK} pipeline as an external post-processing
  step.
\end{itemize}

\subsection{Sister-API design}

The \grip{} package maintains separate APIs for the combinatorial and
weighted engines rather than adding a mode switch to a single function.
This design ensures backward compatibility, makes the
combinatorial/weighted distinction explicit in user code, and allows each
engine to have its own validated parameter presets.  The parameter spaces
are aligned: shared parameters have the same names and semantics, and
weighted-specific parameters (normalization, metric neighbor cap) are
additional rather than replacing existing ones.

\subsection{Speed--quality tradeoffs}

The GRIP-family methods occupy a different region of the
speed--quality Pareto frontier than the KK-family methods.  On small
graphs ($n < 200$), the KK-family is competitive in both speed and
quality because the $\bigO(n^2)$ KK step is cheap at small~$n$.  On
larger graphs, the near-linear scaling of the GRIP multiscale engine
becomes the dominant advantage.

The in-core LGKK integration is particularly efficient because it
amortizes the geodesic cache construction across the multiscale
hierarchy: caches are built incrementally as the active set grows, rather
than being constructed from scratch on the full graph.

\subsection{3D as primary evaluation space}

Our benchmark consistently shows that 3D layouts achieve substantially
lower geodesic error than 2D layouts on the same weighted graph families.
This is not a limitation of any particular algorithm but a consequence of
the intrinsic dimensionality of the test graphs.  For graphs derived from
surfaces, meshes, and 3D bodies, the natural ambient dimension is three,
and forcing the layout into two dimensions introduces irreducible
distortion.

This suggests that researchers working with weighted geometric graphs
should consider 3D layout as the default evaluation space, with 2D used
only when the application specifically requires planar display.

\subsection{Limitations}

The current weighted GRIP implementation does not exploit GPU parallelism,
limiting scalability on very large graphs.  The Dijkstra-based weighted
neighborhood construction is more expensive than BFS-based combinatorial
neighborhoods, though the early-stopping and scratch-reuse optimizations
in Phase~4 of the implementation substantially reduce this cost.  The
benchmark panel covers only small to medium graphs (36--104 vertices); a
broader evaluation at larger scales is an important direction for future
work.

The weighted GRIP engine can produce higher geodesic error than
combinatorial GRIP on some families when used without LGKK refinement,
because the weighted coarsening changes the optimization landscape in
ways that do not always reduce the all-pairs geodesic stress.  LGKK
refinement reliably corrects this, but users should be aware that
weighted GRIP alone is not guaranteed to improve over combinatorial GRIP
on every metric.


% ============================================================
% 10. CONCLUSION
% ============================================================
\section{Conclusion}\label{sec:conclusion}

We have presented a family of extensions that transform GRIP from a
purely combinatorial multiscale graph drawing algorithm into a
geometry-aware multiscale family.  The key contributions are:

\begin{enumerate}
\item A \textbf{weighted multiscale hierarchy} based on weighted MIS
  filtration using Dijkstra shortest-path radii, replacing hop-count
  filtration.
\item \textbf{Weighted neighborhood caches, insertion, and refinement}
  driven by geodesic distances throughout the pipeline.
\item \textbf{Coarse active-set global repulsion} that reduces multiscale
  foldovers in both the combinatorial and weighted engines.
\item \textbf{In-core and post-layout LGKK integration} that provides
  semi-global geodesic refinement within the multiscale framework.
\item A \textbf{systematic benchmark framework} over weighted synthetic
  graph families in 3D.
\end{enumerate}

On a benchmark panel of five weighted graph families, weighted GRIP with
in-core LGKK reduces the mean geodesic error by 22.6\% over
combinatorial GRIP in 3D, while the post-layout LGKK variant achieves a
37.7\% reduction.  The KK-warm-started geodesic polishing family remains
the accuracy ceiling, but the weighted GRIP variants are substantially
faster and do not require external initialization.

All methods are implemented in the open-source \grip{} R package.  The
sister-API design preserves backward compatibility with the existing
combinatorial engine while providing a parallel path for geometry-aware
graph drawing.  The package, benchmark scripts, and all data used in this
paper are publicly available.


% ============================================================
% ACKNOWLEDGMENTS
% ============================================================
\section*{Acknowledgments}

The GRIP algorithm was originally developed in collaboration with
Stephen G.\ Kobourov and Michael T.\ Goodrich.  Research reported in
this publication was supported by grant number INV-048956 from the Gates
Foundation.


% ============================================================
% BIBLIOGRAPHY
% ============================================================
\bibliographystyle{plainnat}
\bibliography{grip-methods-paper}

\end{document}
